\section{Rounding in multiplier-ideal kernel families}
\label{sec:ideal-rounding-kernels}
\index{Bergman kernel!multiplier-ideal rounding}

An intrinsic growth exponent does not by itself determine a kernel
measure limit. At a critical logarithmic coefficient, multiplier-ideal
rounding can leave different leading section counts in even and odd
degrees. The potentials can nevertheless converge to the actual ideal
envelope. The following product makes the growth exponent positive.

Fix $q\geq1$ and put $X=\mathbb P^1\times Y$, where
$Y=(\mathbb P^1)^q$. Write $p_0$ for the first projection.
Let $\sigma_j$ be the Fubini--Study form on
factor $j$, normalized to have integral one, with $j=0$ for the first
factor. Use
\[
 L=\mathcal O(1,\ldots,1),\qquad
 \theta=c_1(L,h)=\sigma_0+\theta_Y,\qquad
 \theta_Y=\sum_{j=1}^q\sigma_j,
\]
where $h$ is the product metric and pullbacks are understood. Sample
with the smooth positive probability
$\nu=\sigma_0\otimes\nu_Y$, where
$\nu_Y=\sigma_1\wedge\cdots\wedge\sigma_q=\theta_Y^q/q!$.
The continuous norm weight is zero.

On the first affine line define
\[
 \psi_0(z)=\tfrac12\log|z|^2-\log(1+|z|^2),\qquad
 \psi=\psi_0\circ p_0.
\]
The inverted chart gives the same expression at infinity. Its canonical
value is $-\infty$ at both endpoints, and
\begin{equation}\label{eq:critical-selector-curvature}
 \theta+\ddc\psi=\tfrac12[D_0]+\tfrac12[D_\infty]+\theta_Y\geq0,
 \qquad D_0=\{0\}\times Y,\quad D_\infty=\{\infty\}\times Y.
\end{equation}
Select $V_m=H^0(X,L^m\otimes\mathcal I(m\psi))$, with its genuine
norm $\int_X|s|_{h^m}^2\,\mathrm d\nu$. Write
$\mathcal B_m=\sum_j|s_{m,j}|_{h^m}^2$ and
$U_m=m^{-1}\log\mathcal B_m$ for an orthonormal basis. Let $\varsigma$
be angular Haar probability on $\{|z|=1\}\subset\mathbb P^1$.

\Needspace{22\baselineskip}
\begin{theorem}\label{thm:ideal-rounding-kernel-limits}
Put $a_m=1$ for even $m$ and $a_m=2$ for odd $m$. Then
\[
 N_m=\dim V_m=a_m(m+1)^q,\qquad
 \lim_{m\to\infty}\frac{\log N_m}{\log m}=q.
\]
The genuine ideal envelope with obstacle zero on $X$ is
\begin{equation}\label{eq:critical-ideal-envelope}
 E_{\theta,X}[\psi]_{\mathcal I}(0)=E=\psi+\log2.
\end{equation}
The norms satisfy Bernstein--Markov, and $U_m\to E$ locally uniformly
off $D_0\cup D_\infty$, in every finite $L^p$ against smooth volume,
and weakly for the closed curvature currents. However,
\begin{align*}
 \frac{q!}{(2k)^q}\mathcal B_{2k}\nu
       &\rightharpoonup q!\,\varsigma\otimes\nu_Y,\\
 \frac{q!}{(2k+1)^q}\mathcal B_{2k+1}\nu
       &\rightharpoonup 2q!\,\varsigma\otimes\nu_Y.
\end{align*}
Thus the kernel measures normalized by $q!/m^q$ do not converge along
all degrees. Dividing instead by the actual dimension gives
$N_m^{-1}\mathcal B_m\nu\rightharpoonup\varsigma\otimes\nu_Y$.
The envelope has zero nonpluripolar top mass:
$\langle(\theta+\ddc E)^{q+1}\rangle=0$.
\end{theorem}
\begin{proof}
By \cref{lem:smooth-center-multiplier-ideal},
\[
 \mathcal I(m\psi)=
 \mathcal O_X(-r_mD_0-r_mD_\infty),\qquad r_m=\lfloor m/2\rfloor.
\]
Global sections in the product affine frame are polynomials of degree
at most $m$ in each variable. Their first exponent satisfies
$r_m\leq\alpha_0\leq m-r_m$, and every other exponent ranges from
$0$ to $m$. There is one first-factor monomial when $m$ is even and
two when it is odd. This proves the dimension formula and shows that
the selection is exactly the stated global multiplier-ideal condition.

To identify the envelope, take any genuine competitor
$v\preceq_{\mathcal I}\psi$ with $v\leq0$ quasi-everywhere.
\Cref{thm:local-i-inputs} gives point Lelong numbers at least $1/2$
on both divisors. Restrict to a first-factor fiber. The restriction is
either identically $-\infty$ or $\sigma_0$-psh. Its endpoint Lelong
numbers are at least the ambient ones: the supremum over a restricted
coordinate ball is at most the ambient supremum, and division by
$\log r^2<0$ reverses the inequality. By the atom identity in
\cref{lem:planar-multiplier-ideal}, its positive curvature measure has
atoms of mass at least $1/2$ at each endpoint. Stokes makes its total
mass one, so these atoms exhaust the measure. The difference from
$\psi_0$ is therefore harmonic on the compact projective line and
constant.

It follows that $v(z,y)=\psi_0(z)+c(y)$ on the complement of the
divisors, allowing $c(y)=-\infty$ on identically singular fibers.
Restriction to $z=1$ gives a genuine $\theta_Y$-psh function
$c(y)=v(1,y)+\log2$, not identically $-\infty$. Shrinking-ball
submeans upgrade the obstacle to $v\leq0$ everywhere. Since
$\max\psi_0=-\log2$, we get $c\leq\log2$ and $v\leq E$.
On the divisors both sides have their canonical value $-\infty$.
Conversely $E\leq0$ and differs from $\psi$ by a constant, so is an
admissible competitor. This proves \eqref{eq:critical-ideal-envelope}.

The one-factor monomial norms are elementary beta integrals. With
$x=|z|^2/(1+|z|^2)$, they are
\begin{equation}\label{eq:fs-monomial-beta-norm}
 \int_{\mathbb P^1}|z^\alpha e^m|_{h^m}^2\,\sigma_0
 =\int_0^1x^\alpha(1-x)^{m-\alpha}\,\mathrm dx
 =\frac{\alpha!(m-\alpha)!}{(m+1)!}.
\end{equation}
Angular orthogonality and product integration diagonalize the global
norm. Summing the full monomial basis in one factor gives the constant
kernel $m+1$ by the binomial formula. Therefore the full kernel on $Y$
is $(m+1)^q$. With $C_k=(2k+1)\binom{2k}{k}$, the actual selected
kernels are
\begin{align}
 \mathcal B_{2k}&=(2k+1)^q C_k\mathrm e^{2k\psi_0},\notag\\
 \mathcal B_{2k+1}&=2(2k+2)^q C_k\mathrm e^{2k\psi_0}.
 \label{eq:critical-even-odd-kernels}
\end{align}
The even formula is used for $k\geq1$, and the odd formula holds also
at $k=0$. In the odd case the two first-factor norms are equal and
their sum contributes the factor $1+|z|^2$, giving the displayed
identity. The full kernel on $X$ is $(m+1)^{q+1}$, so its evaluation
bound gives
$\sup_X|s|_{h^m}\leq(m+1)^{(q+1)/2}\|s\|_m$ for every selected
section. This proves Bernstein--Markov.

The central binomial coefficient is the largest in its row. Hence
$4^k\leq C_k\leq(2k+1)4^k$ and
$\log C_k=2k\log2+O(\log(k+1))$.
Taking logarithms in \eqref{eq:critical-even-odd-kernels} gives
\[
 |U_m-E|\leq\frac{C_q}{m}
       \bigl(1+|\log|z|^2|+\log m\bigr)
 \quad\text{off }D_0\cup D_\infty.
\]
Indeed, $|\psi_0|=\log(2\cosh(\tfrac12\log|z|^2))$ is bounded
by $\log2+\tfrac12|\log|z|^2|$. Logarithmic coordinate poles
have every finite $L^p$ norm against smooth volume. This proves the
potential convergence, and differentiating the $L^1$ limit proves
current convergence. For every $m\geq2$, both canonical potentials
are $-\infty$ on the two divisors.

Let $P_k$ be the probability on the first projective line with density
$C_k\mathrm e^{2k\psi_0}$ against $\sigma_0$.
Its angle is uniform, and $x$ has the beta law with parameters
$(k+1,k+1)$. Its mean is $1/2$ and its variance is
$1/[4(2k+3)]$. Chebyshev's inequality and continuous angular test
functions give $P_k\rightharpoonup\varsigma$.
The exact measure identities are
\[
 \mathcal B_{2k}\nu=(2k+1)^q P_k\otimes\nu_Y,
 \qquad
 \mathcal B_{2k+1}\nu=2(2k+2)^q P_k\otimes\nu_Y.
\]
They prove both subsequential limits and the limit after division by
$N_m$. The different total masses already exclude convergence of the
$q$-normalized family. Ambient-degree normalization by
$(q+1)!/m^{q+1}$ instead gives zero in total variation.

Finally, on the complement of the divisors the curvature in
\eqref{eq:critical-selector-curvature} is $\theta_Y$, of rank $q$.
The nonpluripolar extension charges neither divisor. Its top product
is therefore zero.
\end{proof}

This multiplier-ideal family is not multiplicative. In any two odd
degrees, monomials with minimal order on $D_0$ belong to the selected
spaces, while their product misses the next even ideal condition by
one order. The graded Newton series of \cref{sec:newton-selected-kernels}
has a lattice leading term in every sufficiently large degree.
Here the actual ideal envelope is identified, but the leading
term oscillates. A necessary condition for convergence of any finite
$q$-normalized kernel measures is convergence of $q!N_m/m^q$, by
testing the constant function one. The hypotheses requiring positive
nonpluripolar top mass in \cref{thm:fujiki-partial-kernel-limit} exclude
this critical selector.
